\chapter{The Integrator and the Produced Certificate}
\label{chap:integrator-produced-certificate}

The apparatus of the last several chapters could measure, but it measured in the abstract --- it knew
what an energy was and what it would mean for a path to be stationary, without ever being set running on
a particular path to find out. This chapter sets it running. It builds the integrator, the concrete
instrument that reads the Euler--Lagrange residue off a path and returns a number, and it shows that the
instrument does three things the abstract apparatus could only promise. It reads zero exactly at a
solution and nonzero off one, so its output is a genuine detector of the law of motion. It produces its
convergence rather than assuming it, by exhausting configurations until the residual is forced to zero.
And, most tellingly, it produces a certificate of its own honesty simply by running: a procedure that
leaned on something it could not construct could not be run to a definite answer, so the integrator's
running is itself the receipt that nothing but the one standing grant was used.

\section{The integrator reads the residue}
\label{se:the-integrator-reads-the-residue}

The integrator is the law of motion turned into an instrument. Hand it a path and an action and it
computes the Euler--Lagrange residue --- the amount by which the path fails to be stationary --- and
returns that amount as a reading. Its defining behavior is the one a detector must have. On a genuine
solution, a path that actually satisfies the equation, the integrator reads zero: the identity path, the
simplest solution, is proved to read exactly nothing. On a path that has been kicked off a solution, the
integrator reads a definite nonzero number; the construction's own concrete run perturbs the solution and
watches the reading jump from zero to two. Zero on the solution, nonzero off it: the integrator funges
the residue into a single number that is zero precisely when the path obeys the law and positive
precisely when it does not, and that number is the measurement.

The detector behavior is worth seeing as a worked reading, because it is what makes the integrator an
instrument rather than a definition. Take the identity path, where the output equals the input at every
step --- the plainest solution of the simplest action. Run the integrator on it and the residue at every
node cancels: there is no discrepancy between what the path does and what the equation demands, and the
reading is zero, exactly. Now kick the path: nudge it off the solution at one interior point and leave
the rest. The kicked node no longer satisfies the equation, a residue appears there, and the integrator,
summing the discrepancies, returns a definite positive number --- in the construction's own run, two. The
instrument has tanged the perturbation: it selected, out of a path that looks almost like a solution, the
one place the equation fails, and reported its size. That is what a detector is --- a device that reads
zero on the thing it is tuned to and a definite size on any departure from it --- and the integrator is
the construction's first one built to run.

It is worth being clear about what the integrator does and does not settle, because the boundary is the
same one the tenth chapter drew. Handed a particular path, the integrator computes the residue on that
path, exactly and finitely --- this is the direct residual, the part a finite procedure can always
compute. What it does not do, and cannot, is survey all paths at once to certify that a given one is the
global stationary point; that universal is still the oracle's, the one query the integrator cannot
answer by running. The integrator is precisely the computable face of the law of motion: it reads,
concretely, whether the path in front of it satisfies the equation, and it leaves to the standing grant
the only thing reading cannot reach. The experiment under Laplace's name fixes the status of its output.
Relative to a single committed record, each reading appears determined by what came before --- the
integrator's tape looks like the demon's ledger, the future fixed by the past --- but the demon is local,
demoted to exactly the finite record the integrator has actually written, with no reach past the edge of
what it has read.

\section{Weierstrass by instantiation}
\label{se:weierstrass-by-instantiation}

The twelfth chapter proved the residual could be squeezed toward zero by refinement but left the
convergence itself as a flagged certificate --- a bound asserted to exist. This chapter discharges the
flag, and the way it does so is characteristic. It does not import an analytic theorem that the bound
exists; it constructs the convergence by instantiation, exhausting the configurations until the residual
has nowhere left to be nonzero. As the refinement deepens, the configurations the residual could live on
are enumerated and used up, and a quantity that is eventually zero on every configuration tends to zero
in the only sense a finite construction recognizes. This is Weierstrass approximation done by the
finitist's method: not the existence of an approximating sequence asserted from above, but an
approximating sequence built from below, configuration by configuration, until exhaustion forces the
residual down. The generic is approached, once again, by the instantiation of finite conditions rather
than by a leap to a completed limit, and the convergence that was a hypothesis in the twelfth chapter is
a construction here.

What makes this more than a relabeling of the earlier convergence is the source of the zero. In the
twelfth chapter the residual went to zero because a bound was assumed to drive it there --- the limit was
promised by an analytic certificate the construction marked and did not open. Here the zero has a
constructive cause: the configurations on which the residual could be nonzero are finite at each depth
and enumerable, and the refinement walks through them and uses them up, so that past a certain depth
there is simply no admissible configuration left for a nonzero residual to occupy. The residual is not
asserted to vanish; it is left with nowhere to be. This is the difference between approaching a limit one
hopes is there and approaching one whose every alternative has been struck out, and it is the difference
between the twelfth chapter's flagged certificate and this chapter's produced one. Exhaustion is a thing
a finite construction can do; promising a limit is not, and the integrator does the first in place of
the second.

\section{The boundary is a build artifact}
\label{se:the-boundary-is-a-build-artifact}

Two subtleties the integrator exposes are worth marking, because both are places the construction is
careful not to claim too much. The first is that the boundary --- the anchor that earned definiteness
several chapters ago --- is a build artifact, a feature of how the apparatus is assembled rather than a
fact about the world it measures. The construction needs the boundary to make its forms definite, but it
does not pretend the boundary is anything more than the scaffold's edge, and it marks it as such. The
second is sharper: there is no uniform finite support. For any single configuration the integrator needs
only finitely many entries to read it, but there is no one finite set of entries that suffices for
\emph{all} configurations at once --- any finite set, however large, misses some slip it cannot account
for. The order of the quantifiers is therefore forced and cannot be reversed: for each configuration
there is a finite support, but not a finite support for every configuration. This is the oracle's
universal wearing a different coat. The same gap that made stationarity uncomputable --- a property that
must hold across an unbounded range, with no finite witness covering the whole --- reappears here as the
impossibility of a single finite support, and the construction names it rather than papering over it
with a uniformity it has not got.

The forced quantifier order repays a moment's thought, because it is where the construction's finiteness
and the oracle's infinity meet. To say that for each configuration there is a finite support is to make a
promise the construction can keep: hand it a configuration, and it produces, then and there, the finite
set of entries that reads it. To say there is a finite support for every configuration would be to
promise a single finite object adequate to an unbounded range --- a uniform witness --- and that the
construction cannot keep, because any finite support it offers, a reader can name a configuration that
slips past. The two statements differ only in the order of \emph{for each} and \emph{there is}, and that
order is exactly what separates what a finite procedure can deliver from what only an oracle could. The
construction lives entirely on the deliverable side of that line, and it marks the line rather than
stepping over it --- the same discipline, in the same place, the tenth chapter applied to stationarity
itself.

\section{The certificate is produced}
\label{se:the-certificate-is-produced}

Here is the chapter's most distinctive move, and it is one only a constructive apparatus can make. The
integrator does not merely assert its reading; it \emph{runs}, and produces the reading as an output one
can hold. And that it can be run at all is a certificate of something the construction cares about more
than almost anything else: that the computation is choice-free. A procedure that secretly depended on a
non-constructive grant --- a value asserted to exist without being exhibited --- could not be run to a
definite answer, because there would be no exhibited value to compute with; it would name a result
without being able to produce one. The integrator produces one. It runs to a definite number, and its
running is therefore a receipt, a certificate generated by the act of computing, that everything it
touched was built and nothing was merely assumed. The certificate is produced, not promised: where the
earlier chapters reasoned that the construction was constructive, the integrator demonstrates it by
running, and the demonstration is the run.

It is worth being exact about why running is a certificate and not merely a convenience. The construction
is built so that a procedure can be carried to a definite output only if every value it depends on has
been exhibited --- actually constructed, available to be computed with. A procedure that named a value by
asserting its existence, without producing it, would stall the moment it needed that value, because there
would be nothing there to use. So the set of procedures that run to an answer is exactly the set of
procedures that are constructive through and through. The integrator runs to an answer. That fact, by
itself, places it inside the constructive set and certifies that no existence-without-exhibition was
hiding in it. No separate audit is needed; the run is the audit. This is why the construction prizes a
runnable certificate over a proven one where it can get the runnable kind: a proof that something is
constructive is a claim to be checked, while a run is the constructivity happening, the certificate and
the thing certified collapsing into one act.

The experiment on the ideal ledger names this kind of certificate. A produced certificate is one that
the boundary's reconciliation of the bulk actually generates, an output the apparatus emits rather than
a guarantee it requests; the integrator's run is exactly such an output. And a small economy underlies
it: three tags suffice to type every reading the integrator produces, and one action is refused as
inadmissible, so the whole space of readings is enumerated by a finite, depth-indexed rational skeleton
the integrator can walk. The certificate the integrator produces is therefore not only honest but cheap
--- a finite, runnable receipt, generated on demand, that the measurement it accompanies was taken
constructively and that the only thing assumed anywhere was the single standing grant.

\section{What is demonstrated, and what is assumed}
\label{se:demonstrated-assumed-ch18}

The accounting closes the apparatus side of the book, and it is gratifyingly clean.

What is demonstrated is a working, running instrument and its certificate. The integrator is proved to
read zero on a solution and nonzero off one, and it runs concretely to those values. The convergence of
its residual is constructed by exhaustion, not assumed --- Weierstrass by instantiation. The absence of
a uniform finite support is proved, fixing the quantifier order. And the certificate of choice-freeness
is produced by the act of running, a receipt the computation generates of its own constructivity. Each
is a checked or runnable fact about finite objects.

\begin{quote}
\emph{A note on the one assumption.} It is worth saying precisely what the running integrator does and
does not certify. Its run certifies that the computation is choice-free --- that no non-constructive
grant was smuggled in beneath it. It does \emph{not} certify the one grant the construction has carried
openly since the tenth chapter: that stationarity, the universal over all directions, is decided. That
grant is not a hidden choice the run could expose; it is the declared one, sitting above the
computation, supplying the only thing the integrator's reading cannot reach. The honest summary is that
everything the integrator does is constructive and certified so by its running, on top of exactly one
assumption that no run can discharge because it is not the kind of thing a run touches.
\end{quote}

What is assumed, then, is the standing grant alone, now sharply distinguished from everything the
running integrator certifies away. With a concrete instrument that reads the law of motion, constructs
its own convergence, and produces a receipt of its honesty, the apparatus is complete and proven to run.
The construction has built everything it needs to make a measurement and to vouch for it. What remains is
to point the finished instrument at something and read off a value with a physical name --- which is
what the next chapter does, and the name it reads is the electron.
